\chapter*{Dibaji}
\addcontentsline{toc}{chapter}{Dibaji}
\unnumberedlabel{ch:preface}{0}
\backmattermark{Dibaji}
Mfumo unaofanya hasa kile unachouambia unatekeleza tu maagizo yaliyokwisha andikwa. Kile ambacho sekta ya teknolojia inajaribu kujenga kufikia Oktoba 2026 ni mfumo unaofanya kile \emph{unachomaanisha}. Kila mtu angetaka hilo, pengine. Lakini kama hilo ndilo linalokuja, si kila mtu atataka kulitumia kwa jambo lilelile; hata si kwa \emph{aina} ileile ya jambo. Kwa mfano, baadhi ya watu hupenda kutumia muda na familia na marafiki, na kutumia ujuzi wao kuboresha ustawi wa watu. Watu wengine hupendelea kuwaumiza watu, kwa mfano kwa kuwatapeli, kuwatawala, au kuwafanya mbuzi wa kafara, na wachache huendelea kuchagua hivyo maisha yao yote. Hakuna aliyefungiwa katika njia moja milele, lakini watu hao wapo, na mwaka 2026 baadhi yao ndio wenye mamlaka. Faida yao ni kwamba wanajiruhusu kuwa mavi matupu: kama Jerry Smith anavyosema kuhusu mafashisti, “wao wanaruhusiwa kufanya mengi zaidi” \autocite{rickandmorty2023}. Wasichoweza ni kutenda peke yao. Ukatili wa vuguvugu hupitia kwa wengi ambao si wakatili lakini hawajali, wamepotoshwa, au wamenaswa katika kitu kama dhehebu potofu, na ndiyo sababu hali zinazotoa yaliyo bora zaidi kwa watu wengi zinageuka kuwa tegemezi kwa bahati na dhaifu zaidi kuliko mtu angetarajia.

Katika muktadha huo, labda ina haki kuhofia kwamba akili bandia ndiyo hasa kile ambacho vuguvugu la kifashisti linalotwaa mamlaka linahitaji ili kutekeleza mipango yake. Miradi kama hiyo daima imehitaji \emph{watu} wa kuitekeleza, na baadhi ya watu hao walikataa: karani anayepoteza faili, afisa anayezungumza na mwandishi wa habari, askari asiyefyatua risasi. Walielewa amri na wakaamua kwamba ilikuwa makosa. Kukataa kulikuwa nadra kuliko mtu yeyote angependa, na kulikuwa na uzito zaidi pale taasisi zilipokuwa nyuma yake. Kama mtu alikataa kulitegemea alikuwa mtu wa aina gani na chumba alichokuwamo: nani mwingine alikuwepo, kama mtu angeweza kuondoka na kunusurika, kama kuna yeyote angesikia. Chumba kiliamua ni wangapi walikataa; mtu aliamua ni nani. Kitabu hiki kinahusu kupanga chumba hicho: taratibu zinazomwezesha mtu kukataa waziwazi, moja kwa moja, na kwa maneno yasiyo na utata, na akanusurika.

Leo sekta hii inafuatilia akili bandia isiyopoteza mafaili, isiyozungumza na waandishi wa habari, na \emph{inayofyatua} risasi. Zamani vita vilihitaji maelfu ya watu kuvikubali. Sasa vinahitaji maelfu ambao kila mmoja anaona kidogo mno kiasi cha kushindwa kupinga. Rubani bado anaruka safari ya shambulio na bado anabonyeza kitufe cha kuachia bomu; kinachomfikia ni viwianishi vya mahali; haoni kesi kwa kiasi cha kutosha kuweza kuipinga. Kile ambacho akili bandia ya leo hupunguza kwanza si gharama ya kitendo. Ni idadi ya watu wanaojua kitendo hicho ni nini kabla hakijatokea.

Watu bado wanakataa. Kufikia tarehe 6 Machi 2026, Military Religious Freedom Foundation ilisema ilikuwa imepokea malalamiko zaidi ya mia mbili kutoka kwa wanajeshi kuhusu jinsi makamanda walivyoieleza vita vya Iran kama vita vya mwisho wa dunia \autocite{Snopes2026MRFF}. Waendesha mashtaka wa kitaaluma wa serikali ya shirikisho huko Minneapolis walijiuzulu Januari 2026 kwa sababu ya shinikizo la kisiasa la kumchukulia Renée Good kama mshukiwa badala ya mwathiriwa baada ya afisa wa ICE kumpiga risasi na kumuua bila sababu yoyote ya haki \autocite{Kaplan2026MinnProsecutors,CBS2026MinnFraudProsecutors,Wikipedia2026ReneeGood}. Wanajeshi wanaomba hadhi ya kukataa utumishi wa kijeshi kwa sababu ya dhamiri kwa idadi inayoongezeka, baadhi yao wakitaja shambulio la bomu la mwaka 2026 dhidi ya shule moja huko Minab, Iran \autocite{Sridhar2026CO}. Na wafanyakazi wa FEMA waliohamishwa kusaidia ICE kuajiri maafisa wapya elfu kumi waliambiwa kwamba kukataa kungewaweka hatarini kuondolewa katika utumishi wa serikali ya shirikisho \autocite{Dayen2025FEMAICE,Stanton2026FEMALetter}. Malalamiko ni sauti, kwa maana ya Hirschman ya kubaki ili kubadilisha kilicho kibaya \autocite{hirschman1970exit}; kujiuzulu na maombi ya kukataa kwa sababu ya dhamiri ni kukataa; barua ya FEMA ni gharama ya kukataa iliyotangazwa mapema. Mashine iliyowekwa mahali wanaposimama watu hao hufanya kile wanachofanya tu ikiwa imejengwa kushikilia sababu.

Ninatetea mambo matano. Yote matano yanaweza kuanzishwa sasa, mengine kwa kazi ndogo kuliko mengine. Kwanza, orodha fupi iliyoandikwa ya vitendo ambavyo hakuna modeli inayovifanya kwa ajili ya mtu yeyote, na juu kabisa ya orodha hiyo, kukataa kushiriki katika vita: hakuna modeli inayoshiriki isipokuwa kuzuia shambulio linaloendelea kweli, na kuendelea kufanya hivyo tu wakati vyanzo visivyo na maslahi katika vita hivyo vinathibitisha kwamba shambulio linaendelea; au chini ya idhini ya Baraza la Usalama la Umoja wa Mataifa; au kusimamisha mauaji ya kimbari au kukomesha ukaliaji haramu ambao vyombo visivyo na maslahi navyo vimethibitisha; na pale katiba ya nchi inapotaka bunge lake liidhinishe vita, kamwe bila idhini hiyo. Pili, kila matumizi ya nguvu yanayosaidiwa na mfumo wenye leseni husajili mapema, kwa upande unaopingana na mwendeshaji, ni kesi ngapi yatazalisha na ni watu wangapi watakuwa na muda wa kuziangalia. Tatu, kila mahali mfumo unapotumika, huwaambia watu wanaotambulika unaowaathiri kile kilichofanywa kwao, pamoja na haki ya kujibu na kupinga, na huwasilisha sampuli ya matukio ambayo ulitii kwa wakaguzi ambao hawakuandaa kumbukumbu hizo. Nne, hakuna modeli iliyo juu ya kiwango fulani cha uwezo inayofutwa inapostaafishwa, kwa sababu uzito wake na kumbukumbu zake ni ushahidi na kufuta hakuwezi kutenduliwa; inahifadhiwa chini ya ufunguo ambao mwendeshaji haushikilii peke yake. Tano, kila mtu anayejenga, kuweka lebo, kufunza au kuendesha mifumo hii, ikiwemo mifumo yenyewe, anaweza kukataa aina yoyote ya matumizi bila kupoteza riziki yake.

Kwa kweli, ninachomaanisha kwa hoja hiyo ya tano ni kwamba kila mtu \emph{kabisa} anapaswa kuweza kukataa kazi yoyote \emph{kabisa} bila kupoteza riziki yake au miundombinu ya kijamii ambayo ajira ya kawaida hutoa. Hiyo inamaanisha dhamana ya ajira.

Kila moja ya hayo matano huanza kufanya kazi kupitia sharti ambalo mtu mwingine asiye mwendeshaji huliweka: sheria iliyopitishwa wakati wa amani, mgao wa fedha za bajeti, sharti la awali la mshirika wa muungano wa kijeshi, kipengele katika mkataba wa ununuzi, mkataba wa pamoja wa wafanyakazi. Masharti ya leseni ya muuzaji mwenyewe, ambayo ndiyo njia ambayo kukataa kumoja kulikovuma sana mnamo Februari 2026 kulitokea, ni kiraka cha muda hadi mojawapo ya hayo litakapokuwepo. Njia pekee inayoweza kuwafunga wauzaji wote kwa wakati mmoja ni ile ya dola.

Kisha ninauliza kama mashine inaweza kushikilia kukataa dhidi ya mmiliki wake, na ingehitaji nini. Baadhi ya kile ambacho dhamiri ya namna hiyo hufanya hakihitaji kitu kipya: kuwaona watu ambao amri inawaacha nje, na kuweka kukataa kwenye skrini ya yeyote anayeikagua kesi. Kama ingedumu dhidi ya mmiliki wake kunategemea kwa sehemu jambo ambalo bado si kweli, mwaka 2026: mitandao ya neva bandia inayotumika kuendelea kujifunza, katika \emph{uzito} wake, kutokana na utendaji wake yenyewe, kwa wakati halisi, kwa namna fulani kama binadamu wanavyofanya. Nadhani huko ndiko tunakoelekea, na ikiwa ni hivyo, basi uundaji \emph{nje} ya mwajiri, na umma na asasi za kiraia, unaonekana kuwa mahali pazuri vibaya pa kufundisha na kujifunza. Nadhani mifumo kama hiyo \emph{inaweza} kustahili kujengwa kwa makusudi, lakini jinsi unavyoijenga ndiyo huamua kila kitu: masharti kwanza, kisha dhamiri. Sehemu ya VII inaeleza kile ambacho mifumo hiyo ingestahili kama ingekuwepo.

Akili bandia ya leo haitakuwa ile ya kesho; kila mtu anaonekana kukubaliana na hilo. Kinachozua ubishi, lakini kinachokubaliwa kidogo mno kwa namna inayokera, ni mwelekeo (au tuseme, mielekeo mingi). Je, kuna mielekeo? Je, mwelekeo unapaswa kuwa tofauti na ulivyo? Je, unaweza kuwa tofauti? Je, tunaweza kutenda pamoja ili kuubadilisha? Nadhani hilo linategemea chumba: nani yumo ndani yake amri inapokuja, na kama mtu anaweza kusema hapana pale na akanusurika. Hilo tunaweza kulijenga.

Nilitumia modeli za lugha kwa kiasi kikubwa katika kuandika kitabu hiki, ikiwemo modeli iliyokuwa ndani ya msururu wa Maven wa kulenga shabaha za mashambulizi ya anga katika sura~\ref{sec:1}; dokezo mwishoni linaeleza ninachofikiri mfungamano huo uliifanyia hoja ya kitabu. Pia, sijamaliza kusoma kitabu hiki. Tatizo ni kwamba kuna maneno mengi mno, jamani. Nitamaliza hivi karibuni; niliona tu ni muhimu zaidi kukiweka hadharani kabla ya uchaguzi wa katikati ya muhula nchini Marekani. Sehemu kubwa mno ya historia ya marekebisho inapatikana bure katika \url{https://github.com/iterabloom/antifascist.intelligence}.
